
\documentclass{article} % simplest latex class that gives nice default structures.  Other classes include standalone, book

% Recommended packages:
\usepackage{listings} % used for listing code

\usepackage{graphicx} % used for importing graphics


% Suggested packages:
\usepackage{tikz} % used for drawing figures in LaTeX directly
\usetikzlibrary{automata,calc,arrows} % some useful default tikz packages

\usepackage{enumerate} % used for finer control over enumerate lists

\usepackage{algorithm} % used for specifying algorithms
\usepackage{algpseudocode} % ditto (you need both)
\usepackage{amsmath}
\usepackage{fullpage} % reduces a page's margins to 1cm instead of 1in

\usepackage{mathpazo} % better? fonts

\usepackage{hyperref} % hyperlinks: compile with pdflatex for best results

\usepackage{amssymb}

\begin{document}

\title{COMP9020 Assignment 1}
\author{Tianyi Zhu} % let's keep things anonymous
\date{} % to suppress the date
\maketitle % to actually insert the above three items

\section*{Problem 1} % use section* to suppress numbering
\subsection*{(a)} % use subsection* to suppress numbering
\[\begin{array}{rclr}
S_{2,-4}=\{2\times m+(-4)\times n:m,n\in \mathbb{Z}\}\\
\\
2\times 1+(-4)\times 1&=&-2\\
2\times 1+(-4)\times 2&=&-6\\
2\times 2+(-4)\times 1&=&0\\
2\times 2+(-4)\times 2&=&-4\\
2\times 4+(-4)\times 1&=&4\\
\\
\end{array}\]
-2, -6, 0, -4, 4 are five elements of $S_{2,-4}$.
\subsection*{(b)}
\[\begin{array}{rclr}
S_{12,18}=\{12\times m+18\times n:m,n\in \mathbb{Z}\}\\
\\
12\times 1+18\times 1&=&30\\
12\times 2+18\times 3&=&78\\
12\times 3+18\times 2&=&72\\
12\times 1+18\times 5&=&102\\
12\times 2+18\times 6&=&132\\
\\
\end{array}\]
30, 78, 72, 102, 132 are five elements of $S_{12,18}$.

\subsection*{(c)}
\subsubsection*{(i)}
For $d = \gcd (x,y)$ if and only if, $d\mid x$ and $d\mid y$.\\
Let $x=a\cdot d$ and $y=b\cdot d$ for $a,b\in \mathbb{Z}$.\\
Assume $t\in S_{x,y}$, $t=xe+yf $ for $e,f\in \mathbb{Z}$.\\
\[\begin{array}{rclr}
t&=&xe+yf\\
&=&ade+bdf\\
&=&(ae+bf)d\\
&\because& a,b,e,f\in \mathbb{Z}\\
&\therefore& d\mid t\\
\end{array}\]

Therefore $\forall t\in S_(x,y)$, $d \mid t$ hold. And$S_(x,y)\subseteq \{n:n\in \mathbb{Z}$ and $d\mid n\} $\\
\subsubsection*{(ii)}
Since $z\in S_(x,y)$, let $z=m'x+n'y$ for $m',n'\in \mathbb{Z}$.\\
\[\begin{array}{rclr}
z&=&m'x+n'y\\
&=&m'ed+n'fd\\
&=&(m'e+n'f)d\\
\end{array}\]
\\
$\because z$ is the smallest positive number in $S_{x,y}$ and $d$ is always positive.\\
$\therefore m'e+n'f> 0$\\
$\because m',e,n',f \in \mathbb{Z}$ and $m'e+n'f> 0$\\
$\therefore m'e+n'f\geq 1$\\
\[\begin{array}{rclr}
m'e+n'f&\geq& 1\\
(m'e+n'f)d&\geq& d\\
z&\geq& d\\
d&\leq& z\\
\end{array}\]
\subsection*{d}
\subsubsection*{(i)}
According to the definition of gcd, $\gcd(\vert x\vert,\vert y\vert)$ is equal to $\gcd(x,y).$\\
$\because d = \gcd(\vert x\vert,\vert y\vert)$ is calculatable with Euclid's $\gcd$ Algorithm. Each step in Euclid's $\gcd$ Algorithm can be considered as addition or subtraction.\\
$\therefore$ The result of Euclid's $\gcd$ Algorithm $d$ can be represented in the form of $p\vert x\vert+q\vert y\vert$ where $p,q \in \mathbb{Z}$.\\
$\because$ if x is negative, there always exist a $-m\in \mathbb{Z}$ such that $-mx = m\vert x\vert$.\\
$\therefore$ $d$ can be represented in the form of $px+qy$ where $p,q \in \mathbb{Z}$.\\
$\because d$ satisfy the properites of elements in $ S_{x,y}$\\
$\therefore d\in S_{x,y}$\\
$\because d\leq z$ and $d,z\in S_{x,y}$\\
$\therefore$ if $d \neq z$, $d$ is the smallest positive number in $S_{x,y}$ while $z$ is not. Hence $d=z$\\
$\because d=z$ and $d$ is $\gcd(x,y)$\\
$\therefore x\mid z$ and $y\mid z$\\
\subsubsection*{(ii)}
$\because z=d$\\
$\therefore z\leq d$\\

\section*{Problem 2}
\subsection*{(a)}
$\because \mathit{w}\in [0,y) \cap \mathbb{N}$	\\
$\therefore y>0$\\
There exists $m,n\in \mathbb{Z}$ such $mx+ny =\gcd (x,y)=1$\\
\[\begin{array}{rlcr}
mx+ny&\overset{mod\ y}{=}mx+0\cdot n\\
1&\overset{mod\ y}{=}mx\\
mx&\overset{mod\ y}{=}1\\
\end{array}\]
Assume there exists some $\mathit{w}\in [0,y) \cap \mathbb{N}$ such that $\mathit{w}x\overset{mod\ y}{=}1$.
\[\begin{array}{rlcr}
\because& mx\overset{mod\ y}{=}\mathit{w}x\mbox{ and }x\overset{mod\ y}{=}x\\
\therefore& m\overset{mod\ y}{=}\mathit{w}\\
								&\mathit{w}=m+ky\quad\mbox{${k\in \mathbb{Z}}$}\\
\end{array}\]
If only one $\mathit{w}$ exists in $[0,y)$, it must be the smallest positive $\mathit{w}$ in $\mathit{w}=m+ky$, ${k\in \mathbb{Z}}$ which is $m\%y=\mathit{w}$\\
\[\begin{array}{rclr}
m\%y&=&\mathit{w}\\
m&=&\lfloor \frac{m}{y} \rfloor \cdot y+\mathit{w}\\
(\lfloor \frac{m}{y} \rfloor \cdot y+\mathit{w})x&\overset{mod\ y}{=}1\\
\lfloor \frac{m}{y} \rfloor xy+\mathit{w}x&\overset{mod\ y}{=}1\\
\mathit{w}x&\overset{mod\ y}{=}1\\
\end{array}\]
$\because$For $\mathit{w}=m\%y$, $\mathit{w}x\overset{mod\ y}{=}1$ always hold and $m\%y\in [0,y) \cap \mathbb{N}$.\\
$\therefore$ There is at least one $\mathit{w}\in [0,y) \cap \mathbb{N}$ such that $\mathit{w}x\overset{mod\ y}{=}1$\\
\subsection*{(b)}
Assume $a$ is a common multiple of $x$ and $y$ and $a$ is not a multiple of $l=lcm(x,y)$.\\
\[\begin{array}{rclr}
a&=&e\pmod{l}&\quad\mbox{$e\in (0,l)$}\\
\end{array}\]
$\because x\mid a$ and $y\mid a$\\
$\therefore x\mid e$ and $y\mid e$\\
$\because e<l$\\
$\therefore e$ is the lcm of $x$ and $y$.\\
Hence a common multiple of $x$ and $y$ must be a multiple of $lcm(c,y)$.\\
\[\begin{array}{rclr}
\gcd (x,y)\cdot  lcm(x,y)&=&\vert x\vert \cdot \vert y\vert\\
lcm(x,y)&=&\vert x\vert \cdot \vert y\vert\\
\end{array}\]
Let $kx=z\cdot \vert x\vert \cdot \vert y\vert$ for $z \in \mathbb{Z}$.\\
If $x$ is positive, then $k=z\cdot \vert y\vert$.\\
If $x$ is negative, then $k=-z\cdot \vert y\vert$.\\
$\because y\mid zy$ and $y\mid (-zy)$.\\
$\therefore y\mid k$\\
\subsection*{(c)}
If there is one $\mathit{w'}\in [0,y) \cap \mathbb{N}$ such that $\mathit{w'}x\overset{mod\ y}{=}1$\\
\[\begin{array}{rclr}
&1&\mathit{w'}x\overset{mod\ y}{=}&1\\
&2&\mathit{w'}x\overset{mod\ y}{=}&2\\
&3&\mathit{w'}x\overset{mod\ y}{=}&3\\
&\mbox{...}&&\mbox{...}\\
&(y-1)&\mathit{w'}x\overset{mod\ y}{=}&y-1\\
&y&\mathit{w'}x\overset{mod\ y}{=}&0\\
&(y+1)&\mathit{w'}x\overset{mod\ y}{=}&1\\
\end{array}\]
The next $\mathit{w}$ is $(y+1)\mathit{w'}$.\\
$\because {w'}\neq (y+1)\mathit{w'}$ and $y>0$\\
$\therefore {w'}\neq0$\\
$\because \mathit{w'}\in[1,y) \cap \mathbb{N}$\\
$\therefore (y+1)\mathit{w'}\in [y+1,y(y+1))\cap \mathbb{N}$\\
$\because ([y+1,y(y+1))\cap \mathbb{N})\cap([0,y) \cap \mathbb{N})=\emptyset$\\
$\therefore$ There is at most one $\mathit{w}\in [0,y) \cap \mathbb{N}$ such that $\mathit{w}x\overset{mod\ y}{=}1$\\

\section*{Problem 3}
\[\begin{array}{rclr}
m\% n &=&p &\quad\mbox{$p \in [0,n)$}\\
m&=&k\cdot n +p &\quad\mbox{$k\in \mathbb{Z}$}\\
&\because& m>n\\
&\therefore& k\geq 1\\
\end{array}\]
Then transform the equation required to be proved.
\[\begin{array}{rclr}
\frac{3}{2}(n+(m\%n))&<&m+n\\
\frac{3}{2}(n+p)&<&m+n\\
\frac{3}{2}&<&\frac{m+n}{n+p}\\
\frac{3}{2}&<&\frac{k\cdot n+p+n}{n+p}\\
\frac{3}{2}&<&\frac{k\cdot n}{n+p}+1\\
\frac{1}{2}&<&\frac{k\cdot n}{n+p}\\
\frac{1}{2}&<&k\cdot \frac{n}{n+p}\\
\end{array}\]
It is obvious that the larger $p$ is, the smaller $k\cdot \frac{n}{n+p} is$.
\[\begin{array}{rclr}
\lim\limits_{p\to n}k\cdot \frac{n}{n+p}&<&k\cdot \frac{n}{n+p}\\
k\cdot \frac{n}{n+n}&<&k\cdot \frac{n}{n+p}\\
k\cdot \frac{n}{2\cdot n}&<&k\cdot \frac{n}{n+p}\\
k\cdot \frac{1}{2}&<&k\cdot \frac{n}{n+p}\\
\end{array}\]
The minimum of $k\cdot \frac{1}{2}$ is $\frac{1}{2}$ where $k=1$\\ 
Therefore, for all $m,n\in \mathbb{N}_{>0}$ with $n \leq m$, $\frac{3}{2}(n+(m\%n))<m+n$\\

\section*{Problem 4}
\subsection*{(a)}
\[\begin{array}{rclr}
	A \oplus  A&=& (A \cap  A^{c} ) \cup  (A^{c}  \cap  A) &\quad\mbox{(Definition of $\oplus$)}\\
	&=& \emptyset  \cup  (A^{c}  \cap  A) &\quad\mbox{(Complement with $\cap$)}\\
	&=& (A^{c}  \cap  A) \cup  \emptyset  &\quad\mbox{(Commutatitivity of $\cup$)}\\
	&=& A^{c}  \cap  A &\quad\mbox{(Identity of $\cup$)}\\
	&=& A \cap  A^{c}  &\quad\mbox{(Commutatitivity of $\cap$)}\\
	&=& \emptyset  &\quad\mbox{(Complement with $\cap$)}
\end{array}\]
\subsection*{(b)}
\[\begin{array}{rclr}
    A \cup  \mathcal{U} &=& A \cup  (A \cup  A^{c} ) &\quad\mbox{(Complement with $\cup$)}\\
    &=& (A \cup  A) \cup  A^{c}  &\quad\mbox{(Associativity of $\cup$)}\\
    &=& A \cup  A^{c}  &\quad\mbox{(Idempotence of $\cup$)}\\
    &=& \mathcal{U}  &\quad\mbox{(Complement with $\cup$)}
\end{array}\]
\subsection*{(c)}
\[\begin{array}{rclr}
    A \oplus  B&=& (A \cap  B^{c} ) \cup  (A^{c}  \cap  B) &\quad\mbox{(Definition of $\oplus$)}\\
    &=& ((A \cap  B^{c} ) \cup  A^{c} ) \cap  ((A \cap  B^{c} ) \cup  B) &\quad\mbox{(Distributivity of $\cup$ over $\cap$)}\\
    &=& (A^{c}  \cup  (A \cap  B^{c} )) \cap  ((A \cap  B^{c} ) \cup  B) &\quad\mbox{(Commutatitivity of $\cup$)}\\
    &=& ((A^{c}  \cup  A) \cap  (A^{c}  \cup  B^{c} )) \cap  ((A \cap  B^{c} ) \cup  B) &\quad\mbox{(Distributivity of $\cup$ over $\cap$)}\\
    &=& ((A \cup  A^{c} ) \cap  (A^{c}  \cup  B^{c} )) \cap  ((A \cap  B^{c} ) \cup  B) &\quad\mbox{(Commutatitivity of $\cup$)}\\
    &=& (\mathcal{U}  \cap  (A^{c}  \cup  B^{c} )) \cap  ((A \cap  B^{c} ) \cup  B) &\quad\mbox{(Complement with $\cup$)}\\
    &=& ((A^{c}  \cup  B^{c} ) \cap  \mathcal{U} ) \cap  ((A \cap  B^{c} ) \cup  B) &\quad\mbox{(Commutatitivity of $\cap$)}\\
    &=& (A^{c}  \cup  B^{c} ) \cap  ((A \cap  B^{c} ) \cup  B) &\quad\mbox{(Identity of $\cap$)}\\
    &=& (A^{c}  \cup  B^{c} ) \cap  (B \cup  (A \cap  B^{c} )) &\quad\mbox{(Commutatitivity of $\cup$)}\\
    &=& (A^{c}  \cup  B^{c} ) \cap  ((B \cup  A) \cap  (B \cup  B^{c} )) &\quad\mbox{(Distributivity of $\cup$ over $\cap$)}\\
    &=& (A^{c}  \cup  B^{c} ) \cap  ((B \cup  A) \cap  \mathcal{U} ) &\quad\mbox{(Complement with $\cup$)}\\
    &=& (A^{c}  \cup  B^{c} ) \cap  (B \cup  A) &\quad\mbox{(Identity of $\cap$)}\\
    &=& (B \cup  A) \cap  (A^{c}  \cup  B^{c} ) &\quad\mbox{(Commutatitivity of $\cap$)}\\
    &=& (A \cup  B) \cap  (A^{c}  \cup  B^{c} ) &\quad\mbox{(Commutatitivity of $\cup$)}
\end{array}\]
\subsection*{(d)}
According to uniqueness of complement, if$(A^{c}\cap B^{c})\cup (A\cup B)^{cc}=\mathcal{U}$, then $A^{c}\cap B^{c} =(A\cup B)^{c} $
\[\begin{array}{rclr}
(A^{c}\cap B^{c})\cup (A\cup B)^{cc} &=&(A^{c}\cap B^{c})\cup (A\cup B)&\quad\mbox{(Double complement)}\\
&=& (A\cup B)\cup (A^{c}\cap B^{c})&\quad\mbox{(Commutatitivity of $\cup$)}\\
&=&((A\cup B)\cup A^{c})\cap ((A\cup B)\cup B^{c})&\quad\mbox{(Distributivity of $\cup$ over $\cap$)}\\
&=&(A^{c}\cup (A\cup B))\cap ((A\cup B)\cup B^{c})&\quad\mbox{(Commutatitivity of $\cup$)}\\
&=&((A^{c}\cup A)\cup B)\cap ((A\cup B)\cup B^{c})&\quad\mbox{(Associativity of $\cup$)}\\
&=&(\mathcal{U}\cup B)\cap ((A\cup B)\cup B^{c})&\quad\mbox{(Complement with $\cup$)}\\
&=&(\mathcal{U}\cup B)\cap (A\cup (B\cup B^{c}))&\quad\mbox{(Associativity of $\cup$)}\\
&=&(\mathcal{U}\cup B)\cap (A\cup \mathcal{U})&\quad\mbox{(Complement with $\cup$)}\\
&=&(\mathcal{U}\cup B)\cap (\mathcal{U}\cup A)&\quad\mbox{(Commutatitivity of $\cup$)}\\
&=&\mathcal{U}\cup(A \cap B)&\quad\mbox{(Distributivity of $\cup$ over $\cap$)}\\
&=&(A \cap B)\cup \mathcal{U}&\quad\mbox{(Commutatitivity of $\cup$)}\\

&=&\mathcal{U}&\quad\mbox{(Given in Problem 4(b))}\\

&\because& (A^{c}\cap B^{c})\cup (A\cup B)^{cc}=\mathcal{U}\\
&\therefore& (A^{c}\cap B^{c}) =(A\cup B)^{c}\\
\end{array}\]

\section*{Problem 5}
\subsection*{(a)}
Let $X=\{0, 01\}$,$Y=\{0,11\}$\\
0111 is an element in $(X\cup Y)^{*} $ while it is not an element in $X^*\cup Y^*$\\
Therefore, $(X\cup Y)^{*}=X^*\cup Y^*$ is false.\\
\subsection*{(b)}
Assume $X= \{\lambda \}$:\\
\[\begin{array}{rlcr}
\because& \mbox{For all set A,} XA=A\\
\therefore& X(Y \cup Z)=Y\cup Z\\
&(XY)\cup (XZ)=Y\cup Z\\
\end{array}\]
So $X(Y\cup Z)=(XY)\cup (XZ)$ hold for $X= \{\lambda \}$.\\
\\
Assume $X\neq \{\lambda \}$:\\

The concatenation of set can be represented as $mn$ where $m$ is 1st segment of word and $n$ is 2nd segment.\\
\[\begin{array}{rlcr}
X(Y\cup Z)&=\{mn:m\in X \mbox{ and } n\in (Y\cup Z)\}\\
&=\{mn:m\in X \mbox{ and } n\in Y  \mbox {or } n\in Z\}\\
&=\{mn:m\in X \mbox{ and } n\in Y\}\cup\{mn:m\in X \mbox{ and } n\in Z\}\\
&=(XY)\cup (XZ)\\
\end{array}\]
So $X(Y\cup Z)=(XY)\cup (XZ)$ hold for $X \neq \{\lambda \}$.\\
Therefore, $X(Y\cup Z)=(XY)\cup (XZ)$ is true.\\
\subsection*{(c)}
$\because$Each word segment in $X(X)^{*}$ is from $X$.\\
$\therefore X(X)^{*}\subseteq X^{*}$\\
$\because$Let the 1st word segment pick from $X$ always be $\lambda$. Then $X(X)^{*}=X^{*}$\\
$\therefore X^{*}\subseteq X(X)^{*}$\\
$\because X(X)^{*}\subseteq X^{*}$ and $X^{*}\subseteq X(X)^{*}$\\
$\therefore X(X)^{*}=X^{*}$\\
Therefore, $X(X)^{*}=X^{*}$ is true.\\
\section*{Problem 6}
\subsection*{(a)}
$\{0, 1\}^{\{a,b,c\}}$ has 8 elements:\\
\[\begin{array}{rlcr}
&\{(a,0),(b,0),(c,0)\}\\
&\{(a,0),(b,0),(c,1)\}\\
&\{(a,0),(b,1),(c,0)\}\\
&\{(a,0),(b,1),(c,1)\}\\
&\{(a,1),(b,0),(c,0)\}\\
&\{(a,1),(b,0),(c,1)\}\\
&\{(a,1),(b,1),(c,0)\}\\
&\{(a,1),(b,1),(c,1)\}\\
\end{array}\]
\subsection*{(b)}
$Pow(a,b,c)=\{\emptyset,\{a\},\{b\},\{c\},\{a,b\},\{a,c\},\{b,c\},\{a,b,c\}\}$\\
$Pow(a,b,c)$ is the set of all possible relational pre-image of of $\{0\}$ and $\{1\}$ in $\mathit{f}$.
\[\begin{array}{rlcr}
&R^{\gets}(0)=\{a,b,c\}	&R^{\gets}(1)=\emptyset	&\mbox{for }\mathit{f}=\{(a,0),(b,0),(c,0)\}\\
&R^{\gets}(0)=\{a,b\}	&R^{\gets}(1)=\{c\}		&\mbox{for }\mathit{f}=\{(a,0),(b,0),(c,1)\}\\
&R^{\gets}(0)=\{a,c\}	&R^{\gets}(1)=\{b\}		&\mbox{for }\mathit{f}=\{(a,0),(b,1),(c,0)\}\\
&R^{\gets}(0)=\{a\}		&R^{\gets}(1)=\{b,c\}	&\mbox{for }\mathit{f}=\{(a,0),(b,1),(c,1)\}\\
&R^{\gets}(0)=\{b,c\}	&R^{\gets}(1)=\{a\}		&\mbox{for }\mathit{f}=\{(a,1),(b,0),(c,0)\}\\
&R^{\gets}(0)=\{b\}		&R^{\gets}(1)=\{a,c\}	&\mbox{for }\mathit{f}=\{(a,1),(b,0),(c,1)\}\\
&R^{\gets}(0)=\{c\}		&R^{\gets}(1)=\{a,b\}	&\mbox{for }\mathit{f}=\{(a,1),(b,1),(c,0)\}\\
&R^{\gets}(0)=\emptyset	&R^{\gets}(1)=\{a,b,c\}	&\mbox{for }\mathit{f}=\{(a,1),(b,1),(c,1)\}\\
\end{array}\]
$\forall \mathit{f}, R^{\gets}(0),R^{\gets}(1) \subset Pow(a,b,c)$\\
\subsection*{(c)}
$\{\mathit{w}\in\{0,1\}^{*}:\mathrm{length}(\mathit{w})=3\}=\{000,001,010,011,100,101,110,111\}$\\
$\{\mathit{w}\in\{0,1\}^{*}:\mathrm{length}(\mathit{w})=3\}$ is the set of all possible $\mathit{f}(a)\mathit{f}(b)\mathit{f}(c)$.
\[\begin{array}{rlcr}
&\mathit{f}(a)\mathit{f}(b)\mathit{f}(c)=000	&\mbox{for }\mathit{f}=\{(a,0),(b,0),(c,0)\}\\
&\mathit{f}(a)\mathit{f}(b)\mathit{f}(c)=001	&\mbox{for }\mathit{f}=\{(a,0),(b,0),(c,1)\}\\
&\mathit{f}(a)\mathit{f}(b)\mathit{f}(c)=010	&\mbox{for }\mathit{f}=\{(a,0),(b,1),(c,0)\}\\
&\mathit{f}(a)\mathit{f}(b)\mathit{f}(c)=011	&\mbox{for }\mathit{f}=\{(a,0),(b,1),(c,1)\}\\
&\mathit{f}(a)\mathit{f}(b)\mathit{f}(c)=100	&\mbox{for }\mathit{f}=\{(a,1),(b,0),(c,0)\}\\
&\mathit{f}(a)\mathit{f}(b)\mathit{f}(c)=101	&\mbox{for }\mathit{f}=\{(a,1),(b,0),(c,1)\}\\
&\mathit{f}(a)\mathit{f}(b)\mathit{f}(c)=110	&\mbox{for }\mathit{f}=\{(a,1),(b,1),(c,0)\}\\
&\mathit{f}(a)\mathit{f}(b)\mathit{f}(c)=111	&\mbox{for }\mathit{f}=\{(a,1),(b,1),(c,1)\}\\
\end{array}\]
\section*{Problem 7}
\[\begin{array}{rlcr}
A^{B\times C}&=\{((b,c),a):a\in A \mbox{ and } b\in B \mbox{ and } c\in C\}\\
\\
(A^{B})^{C}	&=\{(b,a):a\in A \mbox{ and } b\in B\}^{C}\\
			&=\{(c,(b,a)):a\in A \mbox{ and } b\in B \mbox{ and } c\in C\}\\
\end{array}\]
Let function $\mathit{f}((b,c),a)=(c,(b,a))$ and $\mathit{f}^{-1}(c,(b,a))=((b,c),a)$ where $a\in A$, $b \in B$ and $c \in C$.\\
$\because$There exists a function $\mathit{f}$ and its inverse function $\mathit{f}^{-1}$ where $\mathit{f}:A^{B\times C} \to (A^{B})^{C}$ and $\mathit{f}^{-1}:(A^{B})^{C} \to A^{B\times C}$.\\ 
$\therefore$ There is a bijection between $A^{B\times C}$ and $(A^{B})^{C}$.\\

\section*{Problem 8}
\subsection*{(a)}
\[\begin{array}{rlcr}
(R_{1};R_{2});R_{3}	&=\{(a,c):\mbox{ there is }a,b\mbox{ with }(a,b)\in R_{1}\mbox{ and }(b,c)\in R_{2}\};R_{3}\\
					&=\{(a,d):\mbox{ there is }a,c\mbox{ with }(a,c)\in R_{1};R_{2}\mbox{ and }(c,d)\in R_{3} \}\\
					&=\{(a,d):\mbox{ there is }a,b\mbox{ with }(a,b)\in R_{1},(b,c)\in R_{2}\mbox{ and }(c,d)\in R_{3}\}\\
					&=\{(a,d):\mbox{ there is }a,b\mbox{ with }(a,b)\in R_{1}\mbox{ and }(b,d)\in R_{2};R_{3} \}\\
					&=R_{1};\{(b,d):\mbox{ there is }b,c\mbox{ with }(b,c)\in R_{2}\mbox{ and }(c,d)\in R_{3}\}\\
					&=R_{1};(R_{2};R_{3})\\
\end{array}\]
\subsection*{(b)}
\[\begin{array}{rlcr}
\because I;R_{1}&=\{(a,b):\mbox{ there is }a\mbox{ with }(a,a)\in I\mbox{ and }(a,b)\in R_{1}\}\\
				&=\{(a,b):\mbox{ there is }a\mbox{ with }(a,b)\in R_{1}\}\\
\therefore I;R_{1}&=R_{1}\\
\\
\because R_{1};I&=\{(a,b):\mbox{ there is }a\mbox{ with }(a,b)\in R_{1}\mbox{ and }(b,b)\in I\}\\
				&=\{(a,b):\mbox{ there is }a\mbox{ with }(a,b)\in R_{1}\}\\
\therefore R_{1};I&=R_{1}\\
\because R_{1};I&=R_{1}\mbox{ and }I;R_{1}=R_{1}\\
\therefore I;R_{1}&=R_{1};I=R_{1}\\
\end{array}\]
\subsection*{(c)}
Let $(a_{1},b_{1})\in R_{1}$, $(a_{2},b_{2})\in R_{2}$ and $(a,b)\in R_{1}\cup R_{2}$\\
\[\begin{array}{rlcr}
(R_{1}\cup R_{2});R_{3}	&=\{(a,c):\mbox{ there is }a,b\mbox{ with }(a,b)\in R_{1}\cup R_{2}\mbox{ and }(b,c)\in R_{3}\}\\
						
\\
(R_{1};R_{3})\cup (R_{2};R_{3})&=\{(a_{1},c):\mbox{ there is }a_{1},b_{1}\mbox{ with }(a_{1},b_{1})\in R_{1}\mbox{ and }(b_{1},c)\in R_{3}\}\cup\\
						&\{(a_{2},c):\mbox{ there is }a_{2},b_{2}\mbox{ with }(a_{2},b_{2})\in R_{2}\mbox{ and }(b_{2},c)\in R_{3}\}\\

\end{array}\]
$\because\forall(a,b)\in  R_{1}\cup R_{2}$, there exist some $(a_{1},b_{1})\in R_{1}$ or $(a_{2},b_{2})\in R_{2}$ equal to $(a,b)$.\\
$\therefore (R_{1}\cup R_{2});R_{3}\subseteq (R_{1};R_{3})\cup (R_{2};R_{3})$\\
$\because\forall(a_{1},b_{1})\in R_{1}$ or $(a_{2},b_{2})\in R_{2}$, it $\in R_{1}\cup R_{2}$.\\
$\therefore (R_{1};R_{3})\cup (R_{2};R_{3}) \subseteq (R_{1}\cup R_{2});R_{3}$\\
$\because (R_{1};R_{3})\cup (R_{2};R_{3}) \subseteq (R_{1}\cup R_{2});R_{3}$ and $(R_{1}\cup R_{2});R_{3}\subseteq (R_{1};R_{3})\cup (R_{2};R_{3})$\\
$\therefore (R_{1}\cup R_{2});R_{3}=(R_{1};R_{3})\cup (R_{2};R_{3})$\\


\subsection*{(d)}
Let $(a,b)\in R_{1}$, $(b_{2},c_{2})\in R_{2}$, $(b_{3},c_{3})\in R_{3}$ and $(b,c)\in R_{2}\cap R_{3}$\\

\[\begin{array}{rlcr}

(R_{1};R_{2})\cap(R_{1};R_{3})	&=\{(a,c_{2}):\mbox{ there is }a,b_{2}\mbox{ with }(a,b_{2})\in R_{1}\mbox{ and }(b_{2},c_{2})\in R_{2}\}\cap\\
								&\{(a,c_{3}):\mbox{ there is }a,b_{3}\mbox{ with }(a,b_{3})\in R_{1}\mbox{ and }(b_{3},c_{3})\in R_{3}\}\\
								&=\{(a,c_{2})\mbox{ or }(a,c_{3}):\mbox{ there is }a,b \mbox{ with }(a,b)\in R_{1},(b,c_2)\in R_{2}\mbox{ and } (b,c_2)\in R_{3}\mbox{ or }\\
								&(b,c_{3})\in R_{3} \mbox{ and }(b,c_{3})\in R_{2}\}\\
								&=\{(a,c):\mbox{ there is }a,b\mbox{ with }(a,b)\in R_{1},(b,c)\in R_{2}\mbox{ and }(b,c)\in R_{3}\}\\
								&=\{(a,c):\mbox{ there is }a,b\mbox{ with }(a,b)\in R_{1}\mbox{ and }(b,c)\in R_{2}\cap R_{3}\}\\
								&=R_{1};(R_{2}\cap R_{3})\\
\end{array}\]
\end{document}
